\chapter{Estimación}
\label{cap:estimacion}

Aquí es donde el método paga su inversión. Estimar un modelo de agentes
heterogéneos sobre series de tiempo era, hasta hace poco, prohibitivo. Con
jacobianos secuenciales es cuestión de minutos.

\section{Del jacobiano a la representación de medias móviles}

El punto de partida es \eqref{eq:ma}. Si los choques exógenos siguen procesos de
medias móviles independientes,
\begin{equation}
d\tilde{Z}_t = \sum_{s\ge0} dZ_s\,\epsilon_{t-s},
\end{equation}
con $\epsilon_t$ innovaciones normales iid, entonces por equivalencia cierta las
variables endógenas siguen
\begin{equation}
d\tilde{X}_t = \sum_{s=0}^{T-1} dX_s\,\epsilon_{t-s},
\qquad dX = \mathbf{G}\,dZ.
\end{equation}
Es decir, \emph{las funciones impulso-respuesta son los coeficientes de la
representación MA}. No hace falta ningún paso adicional.

\section{Momentos analíticos}

De la representación MA salen las autocovarianzas en forma cerrada:
\begin{equation}
\mathrm{Cov}(d\tilde{X}_t, d\tilde{X}_{t'}) = \sum_{s=0}^{T-1-(t'-t)} (dX_s)(dX_{s+t'-t})'.
\label{eq:autocov}
\end{equation}
Nótese que solo dependen de $t'-t$: el proceso es estacionario. La expresión
\eqref{eq:autocov} es una convolución y se evalúa muy rápido con la transformada
rápida de Fourier.

Esto ya permite estimar por el método de momentos simulados --- pero \emph{sin}
simular, que es la gracia: los momentos poblacionales del modelo se calculan
exactamente en milisegundos.

\section{La verosimilitud}

El enfoque estándar en la literatura DSGE es aplicar el filtro de Kalman a la
representación de espacio de estados. Con modelos de agentes heterogéneos ese
espacio de estados es enorme y el filtro se vuelve prohibitivo.

La alternativa es usar directamente la normalidad conjunta. Sea
\begin{equation}
d\tilde{X}^{\text{obs}}_t = \mathbf{B}\, d\tilde{X}_t + u_t
\end{equation}
el vector de $n_{\text{obs}}$ observables, con $u_t$ error de medición normal iid.
Como todo es combinación lineal de normales, el vector apilado de todas las
observaciones es normal multivariado con matriz de covarianzas
\begin{equation}
\mathbf{V} = \mathbb{I}\otimes\Sigma_u + \mathbf{B}\,\mathrm{Cov}(d\tilde{X})\,\mathbf{B}',
\end{equation}
construida a partir de \eqref{eq:autocov}. La log-verosimilitud es la densidad
normal usual:
\begin{equation}
\mathcal{L} = -\tfrac12 \log\det \mathbf{V}
-\tfrac12\,(d\tilde{X}^{\text{obs}})'\,\mathbf{V}^{-1}(d\tilde{X}^{\text{obs}}),
\end{equation}
que se evalúa con una descomposición de Cholesky.

\paragraph{El costo} La descomposición de Cholesky escala como
$n_{\text{obs}}^3 T_{\text{obs}}^3$, lo que empieza a doler cuando la muestra es
larga. Para muestras grandes conviene la aproximación de Whittle, que se calcula
con transformada de Fourier, o construir un espacio de estados a partir de la
representación MA y filtrar con Kalman --- cuyo costo escala linealmente en
$T_{\text{obs}}$ y además permite hacer inferencia sobre los choques con el
suavizador.

\section{Reutilizar jacobianos: la razón de que esto sea viable}

Una estimación bayesiana evalúa la verosimilitud decenas de miles de veces. Lo
que la vuelve factible es que casi nada hay que recalcular.

\begin{receta}[Qué hay que recalcular en cada evaluación]
\begin{enumerate}[leftmargin=1.4em,itemsep=2pt]
\item \textbf{Parámetros de los procesos de choque} (persistencias,
desviaciones estándar). No cambian $\mathbf{G}$ en absoluto: solo cambian
$d\mathbf{Z}$. Se calcula $\mathbf{G}$ \textbf{una vez} y cada evaluación es
álgebra lineal pura. Esta separación no existe en el enfoque de espacio de
estados, y es una ventaja decisiva del espacio de secuencias.
\item \textbf{Parámetros que no tocan el estado estacionario del problema
individual} (rigidez de precios, costos de ajuste, coeficientes de la regla de
Taylor). Cambian los bloques simples pero no el jacobiano del bloque
heterogéneo, que es lo caro. Se recalcula solo la parte barata.
\item \textbf{Parámetros que sí cambian el estado estacionario individual} (la
tasa de descuento, el proceso de ingreso). Obligan a recalcular estado
estacionario y jacobiano en cada evaluación. Sigue siendo mucho más rápido que
antes, pero es el caso costoso y conviene evitarlo: esos parámetros suelen
calibrarse, no estimarse.
\end{enumerate}
\end{receta}

\begin{idea}
Diseñar la estimación consiste, en buena medida, en empujar la mayor cantidad
posible de parámetros a las categorías 1 y 2. Esa es una decisión de
investigación, no solo de implementación: conviene tomarla antes de escribir el
código, no después.
\end{idea}

\section{Una advertencia sobre identificación}

La facilidad de estimar no resuelve el problema de siempre: que los parámetros
estén identificados. Con modelos de agentes heterogéneos hay un riesgo
específico. Muchos mecanismos distintos pueden producir jacobianos agregados
parecidos --- esa es, después de todo, la lección del ``estadístico
suficiente'' --- y si dos parámetros mueven $\J$ de la misma manera, los datos
agregados no pueden separarlos.

La respuesta natural es combinar series de tiempo agregadas con momentos
microeconómicos. El método lo permite: cualquier momento de la distribución
puede ser un producto del bloque (capítulo \ref{cap:extensiones}), y por lo
tanto puede entrar en el vector de observables con su propio error de muestreo.

\begin{tutesis}
Para tu caso la advertencia es concreta. La exposición al choque de liquidez
$a_k$ y el acceso al interbancario $\tau_k$ entran en la decisión del banco de
forma muy parecida: ambos aumentan el costo esperado del faltante y por lo tanto
el colchón voluntario. Con datos agregados de encaje y crédito es improbable que
puedas separarlos. La forma de identificarlos es con el \emph{corte transversal}:
la volatilidad de depósitos por banco disciplina $a_k$ (tu Hecho Estilizado 2) y
los saldos interbancarios por banco disciplinan $\tau_k$. Eso significa que tu
estrategia de IE2 debería apuntar a momentos transversales, no a series de tiempo
agregadas, y conviene decirlo explícitamente en el Anexo II.
\end{tutesis}
